\documentclass[10pt,a4paper]{amsart}
\usepackage{amsmath,amssymb,amsthm,graphicx}
\usepackage{amsmath}
\usepackage{amssymb}
\usepackage{hyperref}
\usepackage[utf8]{inputenc}
\usepackage{csquotes}
\usepackage[foot]{amsaddr}


\usepackage{times}
\usepackage{amsmath,amsthm,amssymb}
\usepackage{array}
\usepackage{bbm}
\usepackage{xcolor}
\usepackage{verbatim}
\usepackage{url}
%\usepackage{authblk}
\pagestyle{plain}


\def\ve{\varepsilon}
\def\vf{\varphi}
\def\vn{\varnothing}

\def\AA{\mathcal{A}}
\def\AAA{\mathbb{A}}
\def\BB{\mathcal{B}}
\def\CC{\mathcal{C}}
\def\CCC{\mathbb{C}}
\def\DDD{\mathbb{D}}
\def\DD{\mathcal{D}}
\def\FF{\mathcal{F}}
\def\HH{\mathcal{H}}
\def\II{\mathcal{I}}
\def\III{\mathbb{I}}
\def\MM{\mathcal{M}}
\def\NN{\mathcal{N}}
\def\NNN{\mathbb{N}}
\def\PPP{\mathbb{P}}
\def\QQQ{\mathbb{Q}}
\def\RR{\mathcal{R}}
\def\RRR{\mathbb{R}}
\def\SS{\mathcal{S}}
\def\SSS{\mathbb{S}}
\def\UU{\mathcal{U}}
\def\VV{\mathcal{V}}
\def\WW{\mathcal{W}}
\def\ZZ{\mathcal{Z}}
\def\ZZZ{\mathbb{Z}}
\def\PP{\mathcal{P}}
\def\TTT{\mathbb{T}}
\def\EEE{\mathbb{E}}


%%%%  jb abbreviations
\def\bK{{\bf K}}
\def\bL{{\bf L}}
\def\bU{{\bf U}}
\def\aut{\mathop{aut}}
\def\betabar{{\bf\beta}}
\def\Uscr{\mathcal U}
\def\myreduct{\upharpoonright}
\def\abar{{\bf a}}
\def\bbar{\bf b}
\def\cbar{\bf c}
\def\lg{\mathop{\rm lg}}
\def\nmdeg{\mathop{\rm nmdeg}}
\def\icl{\mathop{\rm icl}}
% \sidebar{..} puts .. in comment form
\newcommand{\sidebar}[1]{\vskip10pt\noindent
 \hskip.70truein\vrule width2.0pt\hskip.5em
 \vbox{\hsize= 4truein\noindent\footnotesize\relax #1 }\vskip10pt\noindent}



 %%%%

\def\QQQa{\mathbb{Q}\langle \Bar{\alpha},\alpha_0 \rangle}

\def\cc{\mathfrak{c}}

\def\a{\alpha}
\def\b{\beta}
\def\g{\gamma}
\def\d{\delta}
\def\e{\varepsilon}
\def\z{\zeta}
\def\l{\lambda}
\def\k{\kappa}
\def\o{\theta}
\def\r{\rho}
\def\s{\sigma}
\def\t{\tau}
\def\w{\omega}
\def\D{\Delta}

\def\nmdeg{\operatorname{nmdeg}}
\newcommand{\Fmb}[1]{F_{Mb}\left( #1  \right)}
\newcommand{\indf}[1]{F_{ind}\left( #1  \right)}
\newcommand{\fcb}[1]{F_{cb}\left( #1  \right)}
\newcommand{\ficb}[1]{F_{Icb}\left( #1  \right)}





\DeclareMathOperator{\alg}{alg}
\DeclareMathOperator{\RU}{RU}
\DeclareMathOperator{\RM}{RM}
\DeclareMathOperator{\RD}{RD}
\DeclareMathOperator{\tp}{tp}
\DeclareMathOperator{\acl}{acl}
\DeclareMathOperator{\dcl}{dcl}
\DeclareMathOperator{\Cb}{Cb}
\DeclareMathOperator{\dtr}{\delta-trdeg}


\def\id{\operatorname{id}}
\def\Frac{\operatorname{Frac}}
\def\Const{\operatorname{Const}}
\def\spec{\operatorname{Spec}}
\def\span{\operatorname{span}}
\def\exc{\operatorname{Exc}}
\def\Div{\operatorname{Div}}
\def\cl{\operatorname{cl}}
\def\mer{\operatorname{mer}}
\def\trdeg{\operatorname{trdeg}}
\def\ord{\operatorname{ord}}




\newtheorem{theorem}{Theorem}[section]
\newtheorem{thm}[theorem]{Theorem}
\newtheorem{question}[theorem]{Question}
\newtheorem{lemma}[theorem]{Lemma}
\newtheorem{cor}[theorem]{Corollary}
\newtheorem{prop}[theorem]{Proposition}
\newtheorem{fact}[theorem]{Fact}
\newtheorem{claim}[theorem]{Claim}
\newtheorem{subclaim}[theorem]{Subclaim}
\newtheorem{conj}[theorem]{Conjecture}
\newtheorem{cond}[theorem]{Condition}
\newtheorem{observation}[theorem]{Observation}
\newtheorem{notation}[theorem]{Notation}

\theoremstyle{definition}
\newtheorem{definition}[theorem]{Definition}
\newtheorem{example}[theorem]{Example}
\newtheorem{xca}[theorem]{Exercise}
\newtheorem{ques}[theorem]{Question}
\newtheorem{note}[theorem]{Note}

\theoremstyle{remark}
\newtheorem{remark}[theorem]{Remark}
\newcommand{\m}{\mathbb }
\newcommand{\mc}{\mathcal }
\newcommand{\mf}{\mathfrak }
\newcommand{\mb}{F_{Mb} }
\newcommand{\G}{\mathrm{G} }
\newcommand{\U}{\mathrm{U} }



\def \ta {\tau_{\mathcal{D}/\Delta}}
\def \D {\Delta}
\def \DD {\mathcal D}

\newcommand{\gen}[1]{\left\langle#1\right\rangle}
%\renewcommand{\restriction}{\mathord{\upharpoonright}}

\newcommand{\pd}[2]{\frac{\partial #1}{\partial #2}}
\newcommand{\td}{\text{tr.deg.}}
\newcommand{\pp}{\partial }
\newcommand{\pdtwo}[2]{\frac{\partial^2 #1}{\partial #2^2}}
\newcommand{\od}[2]{\frac{d #1}{d #2}}
\def\Ind{\setbox0=\hbox{$x$}\kern\wd0\hbox to 0pt{\hss$\mid$\hss} \lower.9\ht0\hbox to 0pt{\hss$\smile$\hss}\kern\wd0}
\def\Notind{\setbox0=\hbox{$x$}\kern\wd0\hbox to 0pt{\mathchardef \nn=12854\hss$\nn$\kern1.4\wd0\hss}\hbox to 0pt{\hss$\mid$\hss}\lower.9\ht0 \hbox to 0pt{\hss$\smile$\hss}\kern\wd0}
\def\ind{\mathop{\mathpalette\Ind{}}}
\def\nind{\mathop{\mathpalette\Notind{}}}
\numberwithin{equation}{section}
\newcommand{\is}{^{p^ {-\infty}}}
\newcommand{\QQ}{\mathbb Q}
\newcommand{\fh}{\mathfrak h}
\newcommand{\SL}{\operatorname{SL}}


\title{Simple Homogeneous Structures and Indiscernible Sequence Invariants}
\author[1]{John T. Baldwin}
\author[2]{James Freitag}
\author[3]{Scott Mutchnik}
%\affil[1]{Department of Mathematics, Statistics, and Computer Science, University of Illinois at Chicago}
%%%\affil[2]{Department of Biology, University Y}
%%\author{ John T. Baldwin}
%\author{James Freitag} 
\address{Department of Mathematics, Statistics, and Computer Science, University of Illinois at Chicago }
\email{jbaldwin@uic.edu}
\email{jfreitag@uic.edu}

%\author{Scott Mutchnik}
%\address{Department of Mathematics, Statistics, and Computer Science, University of Illinois at Chicago}
\email{mutchnik@uic.edu}

\thanks{JB was partially supported by Simons travel grant G3535. JF was partially supported by NSF CAREER award number 1945251. SM, who is additionally affiliated with the Institut de mathématiques de Jussieu – Paris Rive Gauche, was supported by the NSF under Grant No. DMS-2303034, and by the European Union’s Horizon 2020 research and innovation programme under the Marie Sklodowska-Curie grant agreement No 101034255.}

\date{\today}
\title{The non-minimality degree of the type of a point outside the predicate is r minus one}
\author{John Baldwin, James Freitag, Scott Mutchnik}
\date{arxiv:2405.08211v3; CC BY 4.0}
\begin{document}
\maketitle
\noindent\textit{Editorial note.} All delivered mathematical text is the author's own, copied from the pinned publisher source sources/logic-fl03-simple-homogeneous-structures-and/pinned-original-source.tex; no statement or proof step was rewritten, reordered or removed.\par\medskip
\noindent\textit{Editorial note.} Preamble provenance: the retained preamble is the source's own layout adapter, if any, plus the author's macro block at lines 3-205 of sources/logic-fl03-simple-homogeneous-structures-and/pinned-original-source.tex (everything between the publisher class metadata and \textbackslash{}begin\{document\}, comments excluded). Nothing was added, no macro was invented, and the publisher document class is deliberately not reproduced; the excerpt serves only as preamble, never as a span coordinate.\par\medskip
\noindent\textit{Editorial note.} Metadata provenance: source\_id and source\_version were taken from sources/logic-fl03-simple-homogeneous-structures-and/license-observation.json; the title from SOURCE.json; the author list from the source metadata record; the source URL from the first non-empty of source\_url, source\_urls[0] and the evidence's fixed\_version\_url/license\_url. The licence evidence file used by this card is sources/logic-fl03-simple-homogeneous-structures-and/license-observation.json, recorded in license.evidence\_path so that the binding target is explicit.\par\medskip
\noindent\textit{Editorial note.} External references are navigation pointers only; every deliberately kept step of the author's argument stays inside the counted spans. The displayed numbers of those pointers are the author's own numbering under the author's \textbackslash{}newtheorem declarations (section.index as declared), computed mechanically by the lane's scanner.\par\medskip
\noindent\textit{Editorial note.} The counted statement span contains the author's own commented-out lines (they are part of the source text as staged); they are LaTeX comments and carry no typeset content.\par\medskip
\noindent\textit{Editorial note.} The counted proof contains no reference to any labelled item and no bibliography citation, so the card needs neither a navigation stub nor a bibliography excerpt.\par\medskip
\noindent\textit{Editorial note.} Assembly history: r001 failed with 'LaTeX Error: Option clash for package graphics' because the author's own preamble loads \textbackslash{}usepackage[dvips]\{graphics,epsfig\} while the wrapper class already loads graphicx without options. r002 splits the author preamble block into two spans and omits that single package line; no mathematical macro, statement, span or stub changed. Same family as the earlier \textbackslash{}smaller clash: a publisher package line that must not be re-injected on top of the wrapper baseline.\par\medskip
\noindent\textit{Allowed standard prerequisites (editorial).}\begin{itemize}
\item The theory T\_\{r,k,mu\} and its predicate P, with the relation R on r-tuples of points outside P. Set up in the construction preceding the statement; used inside the counted proof at lines 2354-2355. Supplies the tuple whose existence witnesses the non-forking chain and the relation controlling when a chain has a forking extension.
\item The action of the automorphism group of the model is i-transitive on the complement of P. Used inside the counted proof at lines 2360-2362. Makes each type p\_i, for i < r, the unique non-forking extension of p, which is what keeps the first r-1 steps non-forking.
\end{itemize}\medskip
\begin{theorem}\label{2,2} For $T_{r,k,\mu}$ the  complete type over $\emptyset$, $p$,
  given by $\neg P(x)$ has $\nmdeg(p) =r-1$.
  % and $F_{ind}(p)  > \nmdeg(p)$.
  % The theory $T$ of $M$
%%Satisfies The Conditions Of
%%Construction~\Ref{Const} And So
%has nmdegree $r$.
 \end{theorem}

\begin{proof} Choose distinct $a_1, a_2, a_3,\ldots a_r$ realizing $p(v)$ and such that $R(a_1, a_2,
a_3, \ldots a_r)$.
%Note that each curve has infinitely many points in $\neg
%P(M)$.
%(Note that for every finite $s$ and sequence of free
%variable ${\bf w}_s$ , the formula $\exists {\bf w}_s R(v,\bar)$ in $p$.)
Let $p_i = \tp(a_i/A_{<i})$. Each $p_i$ for $i < r$ does not fork over the
empty set because  $\mathrm{Aut}(M)$ acts $i$-transitively on $\neg P$ guaranteeing
each $p_i$ is the unique non-forking extension of $p$ to its domain $A_{<i}$.
%the forking extension is $x = a_{i-1}$.
The type $p_r$ forks over $\emptyset$ because there are infinitely many
distinct choices of $a_{r-1}^j$ giving different lines. Any pair of the lines
$\ell_i = \langle a_1, a_2 \dots a_{r-1}, a^j_r \rangle$ are infinite and intersect only in
$a_1, a_2 \dots a_{r-1}$ since any further point of intersection would imply
the lines are identical.
 %But then the formula
%
%For $F_{ind}$, let $\langle a_i:i<\omega$ be indiscernible sequence over the
%empty set realizing $p$. Then either every $r$-tuple realizes $R$ or none
%does. In the first case, we have an indiscernible sequence such that
%$\tp(a_n/a_{<n}$ witnesses $F_{ind}(p) > \nmdeg(p)$.
\end{proof}

\end{document}
